\begin{abstract}
Decoy tool calls change both an LLM Agent's encrypted traffic and the work its services perform. We study how call generation affects task recognition and evaluate inference, goal completion, business-state integrity and cost on the same executions. In 3,000 single-host executions with two fixed-weight planners over MCP and TLS~1.3/HTTP/2, ordinary tool calls follow a model decision, whereas controller-inserted cover bypasses that wait. With supplied call windows, label-free timing removal raises recognition from 75 to 112 and from 80 to 108 of 120 tasks. Waits sampled from training traces defeat this shortcut, but attacks using destinations, ordering, sizes and timing recognize 93--106/120; three of four comparisons with selected size-only attacks survive Holm correction. Without supplied call windows, the prespecified removal pipeline loses its advantage, but its idle-gap events rarely match true calls, so this diagnoses the pipeline rather than showing that withheld boundaries protect; recognition remains above chance. A complete shadow task, which also adds a second business action to every episode, reduces the strongest tested call-window counts to 70/120 and 64/120 at about twice the model tokens. These measurements distinguish an exploitable dispatch shortcut from residual trace context and quantify complete-policy costs. A four-axis design space and joint outcome record connect the traffic results to the service capabilities and business-state authority each policy requires.
\end{abstract}
